\documentclass{article}
\usepackage[margin=0.6in]{geometry}
\usepackage{plex-sans}
\usepackage{enumitem}
\usepackage{float}
\usepackage{chngcntr}
\counterwithin{table}{section}
\usepackage{tikz}

\title{SCADS Software Requirements Specification}
\author{jake baldwin}
\date{\today}

\begin{document}
    \maketitle

    \section{Scope}
    The Software Requirements Specification (SRS) identifies the requirements
    for the Spacecraft Cube Attitude Determination System (SCADS).

        \subsection{Document Overview}
            
    
        \subsection{Related Documents}
            \begin{itemize}[label={}]
                \item The Software Development Plan (SDP)
                \item The System and Software Design Description (SSDD)
                \item The Software Test Plan and Processing Procedures (STP)
            \end{itemize}

    \section{Definitions and Acronyms}

        \begin{table}[H]
        \caption{Definitions and Acronyms}
        \vspace{4pt}
        \centering
        \renewcommand{\arraystretch}{1.4}
        \begin{tabular}{|l|p{14cm}|}
        \hline
        \textbf{Requirement ID} & \textbf{Requirements} \\ \hline
        SRS-F-REQ-001 & Spacecraft Cube Attitude Determination System  \\ \hline
        \end{tabular}
        \end{table}


    \section{Functional Requirements}

        \subsection{Required States}
            There exist required states of the spaceship. The finite state machine
            graph is visually displayed in the figure below, and the requirements
            are formalized in the subsequent section.

            \begin{figure}[H]
                \centering
                \fbox{
                    \usetikzlibrary{automata, positioning, arrows.meta}
                    \begin{tikzpicture}[
                        >={Stealth[length=3mm, width=2mm]},
                        shorten >=1pt,
                        node distance=8cm,
                        on grid,
                        auto,
                        every state/.style={minimum size=2.2cm, align=center, font=\small},
                        every edge/.append style={align=center, font=\footnotesize}
                    ]
                        \node[state, initial] (detumble)                         {Detumble};
                        \node[state]          (safe)     [right=of detumble]     {Safe};
                        \node[state]          (nominal)  [below right=of safe]   {Nominal};
                        \node[state]          (man)      [below left=of nominal] {Maneuvering};
                        \node[state]          (attitude) [left=of man]           {Attitude\\Estimation};

                        \path[->]
                            % Detumble outgoing
                            (detumble) edge [pos=0.5, right]
                                            node {Rate\\Low} (attitude)
                                       edge [pos=0.5, above]
                                            node {Sensor Dropout} (safe)

                            % Attitude outgoing
                            (attitude) edge [bend left=15, pos=0.5, left]
                                            node {High\\Rotation Rate} (detumble)
                                       edge [pos=0.5, sloped, above]
                                            node {Uncertainty Low} (nominal)
                                       edge [bend left=15, pos=0.5, sloped, below]
                                            node {Sensor\\Dropout} (safe)

                            % Nominal outgoing
                            (nominal)  edge [bend left=10, pos=0.519, sloped, above]
                                            node {Uncertainty High} (attitude)
                                       edge [bend left=15, pos=0.7, sloped, below]
                                            node {High\\Rotation Rate} (detumble)
                                       edge [bend right=5, pos=0.5, sloped, below]
                                            node {Maneuver Commanded} (man)
                                       edge [bend right=15, pos=0.7, right]
                                            node {Sensor\\Dropout} (safe)

                            % Maneuvering outgoing
                            (man)      edge [bend right=20, pos=0.5, sloped, below]
                                            node {Maneuver Complete} (nominal)
                                       edge [pos=0.5, sloped, below]
                                            node {Uncertainty High} (attitude)
                                       edge [bend left=20, pos=0.75, sloped, above]
                                            node {High\\Rotation Rate} (detumble)
                                       edge [pos=0.5, left]
                                            node {Sensor\\Dropout} (safe)

                            % Safe outgoing
                            (safe)     edge [bend left=30, pos=0.5, sloped, above]
                                            node {Sensors OK+\\Gnd Cmd+\\Rate High} (detumble)
                                       edge [bend left=20, pos=0.5, sloped, above]
                                            node {Sensors OK+\\Gnd Cmd+\\Uncertainty High} (attitude)
                                       edge [bend right=15, pos=0.5, left]
                                            node {Sensors OK+\\Gnd Cmd} (nominal);
                    \end{tikzpicture}
                }
                \caption{SCADS Finite State Machine}
                \label{fig:scads-fsm}
            \end{figure}

        \subsection{State Transition Requirements}

            The following table provides some common definitions used within
            the State Transition Requirements.

            \begin{table}[H]
            \caption{Common Definitions}
            \vspace{4pt}
            \centering
            \renewcommand{\arraystretch}{1.4}
            \begin{tabular}{|l|p{14cm}|}
            \hline
            \textbf{Phrase} & \textbf{Definition} \\ \hline
            Sensor Fault & Refers to a sensor going offline for 
                SD\_CONSECUTIVE\_NUM\_SAMPLES consecutive samples or 
                SD\_TRIGGER\_AGGREGATE\_SAMPLES samples per the last 
                SD\_AGGREGATE\_SAMPLES. \\ \hline
            Sensor Recover & Refers to a sensor coming back online for 
                SR\_CONSECUTIVE\_NUM\_SAMPLES consecutive samples or 
                SR\_TRIGGER\_AGGREGATE\_SAMPLES samples per the last 
                SR\_AGGREGATE\_SAMPLES. \\ \hline
            High Rotation & When the measured rate of rotation exceeds 
                ROTATION\_RATE\_HIGH radians per second over ROTATION\_SAMPLES\_HIGH 
                consecutive samples. \\ \hline
            Low Rotation & When the measured rate of rotation is less than
                or equal to ROTATION\_RATE\_LOW radians per second over 
                ROTATION\_SAMPLES\_LOW consecutive samples. \\ \hline
            High Uncertainty & When the measured attitude uncertainty exceeds 
                UNCERATINTY\_HIGH units per UNCERTAINTY\_SAMPLES\_HIGH 
                consecutive samples. \\ \hline
            Low Uncertainty & When the measured attitude uncertainty is
                less than UNCERATINTY\_LOW units per UNCERTAINTY\_SAMPLES\_LOW 
                consecutive samples. \\ \hline
            \end{tabular}
            \end{table}

            The following table defines the state transition requirements 
            between the aforementioned states.
            
            \begin{table}[H]
            \caption{State Transition Requirements}
            \vspace{4pt}
            \centering
            \renewcommand{\arraystretch}{1.4}
            \begin{tabular}{|l|p{9cm}|p{5cm}|}
            \hline
            \textbf{Requirement ID} & \textbf{Requirement} & \textbf{Acceptance Criteria} \\ \hline
            SCADS-ST-REQ-001 & The SCADS shall transition from attitude 
                                estimation to nominal when the sensors
                                indicate uncertainty low. 
                                & TBD - STPP \\ \hline
            SCADS-ST-REQ-002 & The SCADS shall transition from nominal 
                                to detumble with the highest precedence if the
                                sensors indicate high rotation.& TBD - STPP \\ \hline
            SCADS-ST-REQ-003 & The SCADS shall transition from nominal 
                                to safe with the second highest precedence if 
                                a sensor faults. & TBD - STPP \\ \hline
            SCADS-ST-REQ-004 & The SCADS shall transition from nominal 
                                to attitude estimation with the third highest 
                                precedence if the sensors indicate high 
                                uncertainty. & TBD - STPP \\ \hline
            SCADS-ST-REQ-005 & The SCADS shall transition from nominal 
                                to maneuvering when a maneuver is commanded. 
                                & TBD - STPP \\ \hline
            SCADS-ST-REQ-006 & The SCADS shall transition from detumble 
                                to attitude estimation when the sensors
                                indicate low rotation. 
                                & TBD - STPP \\ \hline
            SCADS-ST-REQ-007 & The SCADS shall transition from maneuvering
                                to nominal when the commanded maneuver is 
                                completed. & TBD - STPP \\ \hline
            SCADS-ST-REQ-008 & The SCADS shall transition from attitude 
                                estimation to safe if a sensor faults. 
                                & TBD - STPP \\ \hline
            SCADS-ST-REQ-009 & The SCADS shall transition from maneuvering 
                                to detumble when sensors indicate high rotation. 
                                & TBD - STPP \\ \hline
            SCADS-ST-REQ-010 & The SCADS shall transition from detumble 
                                to safe when a sensor faults.
                                & TBD - STPP \\ \hline
            SCADS-ST-REQ-011 & The SCADS shall transition from attitude 
                                estimation to detumble when sensors indicate 
                                high rotation.
                                & TBD - STPP \\ \hline
            SCADS-ST-REQ-012 & The SCADS shall transition from maneuvering to
                                attitude estimation when sensors indicate 
                                uncertainty high.
                                & TBD - STPP \\ \hline
            SCADS-ST-REQ-013 & The SCADS shall transition from maneuvering to
                                safe when a sensor faults.
                                & TBD - STPP \\ \hline
            SCADS-ST-REQ-014 & The SCADS shall transition from safe to
                                detumble with the highest precedence when 
                                sensors recover, a ground command is received, 
                                and sensors indicate high rotation. 
                                & TBD - STPP \\ \hline
            SCADS-ST-REQ-015 & The SCADS shall transition from safe to
                                attitude estimation with the second highest 
                                precedence when sensors recover, a 
                                ground command is received, and sensors 
                                indicate uncertainty high. 
                                & TBD - STPP \\ \hline
            SCADS-ST-REQ-016 & The SCADS shall transition from safe to
                                nominal with the third highest precedence
                                when sensors recover, and a ground command is 
                                received. 
                                & TBD - STPP \\ \hline
            SCADS-ST-REQ-017 & The SCADS shall enter detumble when the system
                                is powered on. & TBD - STPP \\ \hline
            SCADS-ST-REQ-018 & Upon any transition out of maneuvering prior to
                                maneuver completion, SCADS shall command all
                                thrusters to the off state. & TBD - STPP \\ \hline


            \end{tabular}
            \end{table}

        \subsection{Fault Detection Requirements}
            The fault detection requirements are defined in the table below.

            \begin{table}[H]
            \caption{Fault Detection Requirements}
            \vspace{4pt}
            \centering
            \renewcommand{\arraystretch}{1.4}
            \begin{tabular}{|l|p{9cm}|p{5cm}|}
            \hline
            \textbf{Requirement ID} & \textbf{Requirement} & \textbf{Acceptance Criteria} \\ \hline
            SCADS-FD-REQ-001 & The SCADS shall detect a Sensor Fault for each 
                             sensor, where Sensor Fault is defined as a sensor 
                             going offline for SD\_CONSECUTIVE\_NUM\_SAMPLES 
                             consecutive samples or 
                             SD\_TRIGGER\_AGGREGATE\_SAMPLES samples per the 
                             last SD\_AGGREGATE\_SAMPLES. 
                           & TBD - STPP \\ \hline
            SCADS-FD-REQ-002 & The SCADS shall detect a sensor out-of-range fault 
                             for each sensor when any reading exceeds its 
                            specified minimum or maximum valid range 
                            [Sensor Specification, TBD]. 
                           & TBD - STPP \\ \hline
            SCADS-FD-REQ-003 & The SCADS shall detect a task overrun fault when 
                             any critical task fails to complete execution 
                             within its allocated period.
                           & TBD - STPP \\ \hline
            SCADS-FD-REQ-004 & The SCADS shall recognize a watchdog-induced reset 
                             condition upon boot and log it as a watchdog 
                             timeout fault.
                           & TBD - STPP \\ \hline
            SCADS-FD-REQ-005 & The SCADS shall detect an estimator divergence 
                             fault when the attitude uncertainty exceeds 
                             UNCERTAINTY\_DIVERGENCE\_THRESHOLD or increases 
                             monotonically for DIVERGENCE\_SAMPLES consecutive 
                             samples.
                           & TBD - STPP \\ \hline
            SCADS-FD-REQ-006 & Upon detection of any fault defined in 
                             SCADS-FD-REQ-001 through SCADS-FD-REQ-005, 
                             the SCADS shall log the fault type, timestamp, 
                             and originating subsystem to non-volatile storage.
                           & TBD - STPP \\ \hline
            \end{tabular}
            \end{table}

        \subsection{Core Functional Requirements}

            Attitude estimation (run the estimator, output current quaternion + uncertainty)
            Attitude control (compute actuator commands from desired vs. current attitude)
            Sensor data acquisition (read from sensors at some rate, convert raw to physical units)
            Actuator command generation (send commands to actuators at some rate)
            Command processing (accept ground commands, validate, dispatch)
            Telemetry generation (package state data, send to ground)
            Health monitoring (track sensor/actuator/task health, feed into fault detection)
            Time management (maintain mission time, handle time correlation)

            \begin{table}[H]
            \caption{Functional Requirements}
            \vspace{4pt}
            \centering
            \renewcommand{\arraystretch}{1.4}
            \begin{tabular}{|l|p{9cm}|p{5cm}|}
            \hline
            \textbf{Requirement ID} & \textbf{Requirements} & \textbf{Acceptance Criteria} \\ \hline
            SCADS-F-REQ-001 & Spacecraft Cube Attitude Determination System & AC \\ \hline
            \end{tabular}
            \end{table}
        \subsection{Fault Detection Requirements}

    \section{Performance Requirements}

        Control loop rate (Hz)
        Attitude estimation rate (Hz)
        Command response latency (ms)
        Telemetry rate (Hz)
        Boot time to operational (s)
        Pointing accuracy (degrees, if you care)
        Rate damping performance (time to reduce from X rad/s to Y rad/s in Detumble)

        \begin{table}[H]
        \caption{Performance Requirements}
        \vspace{4pt}
        \centering
        \renewcommand{\arraystretch}{1.4}
        \begin{tabular}{|l|p{9cm}|p{5cm}|}
        \hline
        \textbf{Requirement ID} & \textbf{Requirements} & \textbf{Acceptance Criteria} \\ \hline
        SRS-P-REQ-001 & Spacecraft Cube Attitude Determination System & AC \\ \hline
        \end{tabular}
        \end{table}

    \section{Interface Requirements}

        Ground command interface (e.g., "shall accept commands over UART at 115200 baud")
        Telemetry output interface (same or separate UART)
        Sensor input interface (whatever your simulator uses — file I/O, shared memory, UART from a host)
        Actuator output interface (similar, whatever drives your actuator model)

        \begin{table}[H]
        \caption{Interface Requirements}
        \vspace{4pt}
        \centering
        \renewcommand{\arraystretch}{1.4}
        \begin{tabular}{|l|p{9cm}|p{5cm}|}
        \hline
        \textbf{Requirement ID} & \textbf{Requirements} & \textbf{Acceptance Criteria} \\ \hline
        SRS-I-REQ-001 & Spacecraft Cube Attitude Determination System & AC \\ \hline
        \end{tabular}
        \end{table}

    \section{Safety and Fault Tolerance Requirements}
        "Upon fault detection, the ADCS shall transition to Safe Mode within [TBD] ms."
        "The ADCS shall maintain sufficient power margin in Safe Mode to operate for [TBD] hours."
        "The ADCS shall log fault events to non-volatile memory."
        "No single sensor failure shall cause loss of attitude determination" (aspirational — depends on sensor suite)
        \begin{table}[H]
        \caption{Safety and Fault Tolerance Requirements}
        \vspace{4pt}
        \centering
        \renewcommand{\arraystretch}{1.4}
        \begin{tabular}{|l|p{9cm}|p{5cm}|}
        \hline
        \textbf{Requirement ID} & \textbf{Requirements} & \textbf{Acceptance Criteria} \\ \hline
        SRS-SFT-REQ-001 & Spacecraft Cube Attitude Determination System & AC \\ \hline
        \end{tabular}
        \end{table}


    \section{Operational Environment Requirements}
        Orbital parameters (LEO, altitude, inclination — drives thermal, radiation, eclipse timing)
        Radiation environment (SEU rate expectations, if hardened)
        Temperature range
        Duty cycle (continuous operation vs. eclipse-only)
        \begin{table}[H]
        \caption{Operational Environment Requirements}
        \vspace{4pt}
        \centering
        \renewcommand{\arraystretch}{1.4}
        \begin{tabular}{|l|p{9cm}|p{5cm}|}
        \hline
        \textbf{Requirement ID} & \textbf{Requirements} & \textbf{Acceptance Criteria} \\ \hline
        SRS-OE-REQ-001 & Spacecraft Cube Attitude Determination System & AC \\ \hline
        \end{tabular}
        \end{table}

\end{document}

